\documentclass[11pt]{article}
\usepackage[a4paper,margin=2.2cm]{geometry}
\usepackage{xcolor}
\usepackage{fontspec}
\usepackage{xeCJK}
\usepackage{enumitem}
\setmainfont{Times New Roman}
\setCJKmainfont{SimSun}
\setlist[itemize]{leftmargin=1.4em,itemsep=0.35em,topsep=0.35em}
\setlength{\parindent}{0pt}
\setlength{\parskip}{0.55em}
\pagestyle{plain}
\begin{document}
{\color[HTML]{1F5FD2}\LARGE\bfseries 家长培训手册：早期干预的重要性\par}
{\color[HTML]{667085}\small 星轨原创教育资源：用于家庭与学校共同制定支持计划。\par}
\vspace{0.8em}
{\color[HTML]{111827}\large\bfseries 早期干预价值\par}
\begin{itemize}
\item 早期干预通过日常活动促进沟通、游戏、动作、自理和社会参与，重点是让练习自然发生在生活中。
\item 家长不需要变成治疗师，而是作为最了解儿童生活的人，与专业人员共同选择可持续的策略。
\end{itemize}
{\color[HTML]{111827}\large\bfseries 家庭练习\par}
\begin{itemize}
\item 每天选择两个稳定流程，例如点心、穿衣、洗澡或睡前阅读，把目标设计得小而可重复。
\item 先跟随儿童兴趣并模仿，再加入一个新动作或新词，给儿童足够等待时间。
\item 记录有效动机、困难时段和成功策略，便于团队复盘时调整计划。
\end{itemize}
{\color[HTML]{111827}\large\bfseries 照护者支持\par}
\begin{itemize}
\item 干预计划必须可持续。如果某个策略让家庭长期疲惫，应降低频率或与团队重新排序目标。
\item 照护者的休息、情绪和获得支持的渠道，也是儿童发展环境的重要组成部分。
\end{itemize}
\vfill
{\color[HTML]{667085}\footnotesize 本材料为应用内原创教育内容，不能替代医疗、法律或正式评估建议。\par}
\end{document}
